\section{Motivación e introducción}
La segmentación automática de vasos retinianos es una tarea central en sistemas de apoyo al diagnóstico de retinopatía diabética, hipertensión y otras patologías microvasculares. Desde el punto de vista clínico, una máscara vascular precisa permite cuantificar tortuosidad, calibre y densidad de vasos; desde el punto de vista computacional, el problema es desafiante por el fuerte desbalance entre fondo y vasos, la variación de iluminación entre adquisiciones y la presencia de estructuras extremadamente finas.

El objetivo de este trabajo es construir un sistema reproducible, completamente implementado en PyTorch y alineado con los entregables del examen parcial. En particular, el repositorio desarrolla una U-Net desde cero, incorpora variantes con compuertas de atención y \emph{nested skip connections}, compara decisiones de entrenamiento mediante ablación, evalúa el rendimiento dentro de DRIVE y analiza tipos de fallo desde una perspectiva arquitectónica.

Los entregables cubiertos por esta plantilla son los siguientes:
\begin{itemize}
    \item Implementación personalizada de U-Net y variantes sin depender de bibliotecas de segmentación.
    \item Estudio de ablación sobre decisiones arquitectónicas y de pérdida.
    \item Evaluación en DRIVE con sensibilidad, especificidad, F1 y AUC-ROC.
    \item Comparación arquitectónica en DRIVE para medir ganancia relativa de cada variante.
    \item Análisis cualitativo de fallos por tipo de vaso.
    \item Evidencia de estrategias de robustez (preprocesamiento y TTA) dentro del dominio DRIVE.
\end{itemize}
